\section*{Capítulo \ref{ch:b1:structure-spectroscopy} --- La espectroscopia al servicio de la estructura}

\begin{solution}{exo:b1:structure-spectroscopy:1}
$\varepsilon = A/(\ell c) = 0.64/(1.00 \times \num{2.0e-5}) =
\qty{3.2e4}{L/(mol.cm)}$. Con el triple de concentración: $A = 1.92$, si
la ley sigue cumpliéndose a esa \omterm{def:b1:structure-spectroscopy:absorbance}{absorbancia}.
\end{solution}

\begin{solution}{exo:b1:structure-spectroscopy:2}
\ce{C6H6}: $D = (12 + 2 - 6)/2 = 4$, el benceno (un anillo y tres enlaces
$\pi$). \ce{C4H8O2}: 1, el etanoato de etilo o el ácido butanoico.
\ce{C3H7NO}: $(6 + 2 + 1
- 7)/2 = 1$, la propanamida. \ce{C2H3Cl}: 1, el cloroeteno. \ce{C10H14O}:
$(20 + 2 - 14)/2 = 4$: la carvona tiene un anillo, dos C=C y un C=O.
\end{solution}

\begin{solution}{exo:b1:structure-spectroscopy:3}
Propanona: una señal (6H). Ácido etanoico: dos (3H, 1H). 2-Metilpropano:
dos (9H, 1H). 1,2-Dicloroetano: una (4H).
\end{solution}

\begin{solution}{exo:b1:structure-spectroscopy:4}
$\delta = 1460/400 = \qty{3.65}{ppm}$; a \qty{90}{MHz}, la señal está a
$3.65 \times 90 = \qty{329}{Hz}$ de la referencia. Las \omterm{def:b1:structure-spectroscopy:coupling}{constantes de
acoplamiento} no cambian en hercios, mientras que las diferencias de
desplazamiento crecen con el campo: los \omterm{def:b1:structure-spectroscopy:coupling}{multipletes} que se solapan a
campo bajo se separan a campo alto.
\end{solution}

\begin{solution}{exo:b1:structure-spectroscopy:5}
\ce{CH3}: triplete (3H); \ce{CH2} central: cinco vecinos, sextuplete
1\,:\,5\,:\,10\,:\,10\,:\,5\,:\,1 (2H); \ce{CH2Br}: triplete (2H), el más
desapantallado.
\end{solution}

\begin{solution}{exo:b1:structure-spectroscopy:6}
Etanol: O--H (3666) y ningún C=O; propanona: C=O (1738) y ningún O--H;
ácido etanoico: ambas, O--H (3582) y C=O (1788).
\end{solution}

\begin{solution}{exo:b1:structure-spectroscopy:7}
Etanoato de etilo: cuadruplete (2H) hacia 4.1, singlete (3H) hacia 2.0 y
triplete (3H) hacia 1.3 ppm. Propanoato de metilo: singlete (3H) del
\ce{OCH3}, desapantallado por el oxígeno (hacia 3.7 ppm), cuadruplete
(2H) del \ce{CH2C=O} (hacia 2.3) y triplete (3H) hacia 1.1. El
cuadruplete es la señal más desapantallada en el primero y no en el
segundo: un cuadruplete hacia 4 ppm indica \ce{O-CH2CH3}.
\end{solution}

\begin{solution}{exo:b1:structure-spectroscopy:8}
Líneas separadas \qty{0.080}{ppm}, es decir, $0.080 \times 90 = \qty{7.2}{Hz}$:
la \omterm{def:b1:structure-spectroscopy:coupling}{constante de acoplamiento}. El triplete tiene la misma separación: el
\ce{CH2} y el \ce{CH3} están acoplados entre sí.
\end{solution}

\begin{solution}{exo:b1:structure-spectroscopy:9}
Una banda C=O y una sola señal: la propanona, \ce{CH3COCH3}. El isómero
aldehído, el propanal \ce{CH3CH2CHO}, presenta la tensión C--H del grupo
\ce{CHO} y su protón muy desapantallado.
\end{solution}

\begin{solution}{exo:b1:structure-spectroscopy:10}
El primer conjunto desdobla la línea en $n_1 + 1$ líneas; cada una de
ellas la desdobla el segundo conjunto en $n_2 + 1$: $(n_1 + 1)(n_2 + 1)$
líneas. Si $J_1 =
J_2$, la posición de una línea solo depende de $k_1 + k_2$, que toma
$n_1 + n_2 +
1$ valores: la regla $n + 1$ con $n = n_1 + n_2$. \ce{CH2} central del
1-bromopropano: $4 \times 3 = 12$ líneas en general, seis (un sextuplete)
cuando las dos constantes son iguales.
\end{solution}

\begin{solution}{exo:b1:structure-spectroscopy:11}
$D = 0$; un alcohol (banda O--H, sin C=O). Doblete (6H): dos \ce{CH3}
equivalentes sobre un mismo CH; \omterm{def:b1:structure-spectroscopy:coupling}{multiplete} (1H): ese CH; doblete (2H): un
\ce{CH2} contiguo al CH y que lleva el OH; singlete ancho: el OH.
2-Metilpropan-1-ol, \ce{(CH3)2CH-CH2-OH}.
\end{solution}

\begin{solution}{exo:b1:structure-spectroscopy:12}
Cada especie absorbe de forma independiente: $A = A_1 + A_2 = \ell(\varepsilon_1c_1
+ \varepsilon_2c_2)$. A dos longitudes de onda, con los cuatro
coeficientes conocidos, dos ecuaciones lineales dan $c_1$ y $c_2$.
\end{solution}

\begin{solution}{pb:b1:structure-spectroscopy:1}
\textbf{1.} $0.5690 \times 12.0/44.0 = \qty{0.1552}{g}$.
\textbf{2.} $0.2328 \times 2.0/18.0 = \qty{0.02587}{g}$.
\textbf{3.} $0.2500 - 0.1552 - 0.0259 = \qty{0.0689}{g}$.
\textbf{4.} C \qty{62.1}{\percent}, H \qty{10.3}{\percent}, O
\qty{27.6}{\percent}.
\textbf{5.} C $0.1552/12.0 = 0.01293$, H $0.02587$, O $0.0689/16.0 =
0.00431$ mol: razón $3.00 : 6.00 : 1$, \ce{C3H6O}.
\textbf{6.} $M = \rho RT/p = 3740 \times 8.314 \times 373.15/(1.000 \times 10^5) =
\qty{116}{g/mol}$ (densidad en \unit{g/m^3}).
\textbf{7.} $116/58 = 2$: \ce{C6H12O2}, $D = (12 + 2 - 12)/2 = 1$.
\textbf{8.} Una tensión C=O.
\textbf{9.} Cualquier O--H: ni ácido ni alcohol.
\textbf{10.} Una tensión C--O, intensa: el C--O de un éster.
\textbf{11.} Tensiones C--H de grupos \ce{CH2} y \ce{CH3}.
\textbf{12.} Un enlace $\pi$, ocupado por el C=O; con un C--O intenso y
sin O--H, un éster. Un ácido necesita un O--H; un diol no tiene C=O. Una
cetona con un grupo éter tendría su C=O cerca de los \qty{1738}{cm^{-1}}
de la propanona, mientras que los \qty{1754}{cm^{-1}} observados están
cerca de los \qty{1764}{cm^{-1}} del etanoato de etilo; la RMN de la parte
III confirma el éster.
\textbf{13.} Cinco; $2 + 2 + 2 + 3 + 3 = 12$, los doce hidrógenos.
\textbf{14.} Un \ce{CH2} con tres vecinos (un \ce{CH3}), unido al oxígeno
(desapantallado): \ce{-O-CH2-CH3}.
\textbf{15.} El \ce{CH3} de ese grupo etilo, con dos vecinos; las dos
señales están acopladas entre sí, de ahí una sola $J$.
\textbf{16.} Un \ce{CH2} con dos vecinos, contiguo al C=O.
\textbf{17.} Un \ce{CH2} con cinco vecinos: entre un \ce{CH2} y un
\ce{CH3}.
\textbf{18.} Un \ce{CH3} con dos vecinos, alejado del grupo funcional.
\textbf{19.} \ce{CH3CH2O-} y \ce{-CH2CH2CH3}.
\textbf{20.} Los protones de 4.13 ppm están en el carbono unido al
oxígeno, y los de 2.28 ppm, junto al C=O: ambos atractores de electrones.
\textbf{21.} \ce{CH3CH2-O-CO-CH2CH2CH3} o \ce{CH3CH2-CO-O-CH2CH2CH3}. Solo
el primero sitúa sobre el oxígeno el \ce{CH2} con tres vecinos (el
cuadruplete de 4.13 ppm).
\textbf{22.} El propanoato de propilo es el segundo: su \ce{OCH2} sería un
triplete y el cuadruplete estaría hacia 2.3 ppm.
\textbf{23.} El etanoato de butilo, \ce{CH3CO-O-C4H9}, daría un singlete
(3H) hacia 2.0 ppm y ningún cuadruplete.
\textbf{24.} El butanoato de etilo, \ce{CH3CH2CH2-CO-O-CH2CH3}.
\textbf{25.} IR: C=O 1754, C--O 1182, C--H 2982, sin O--H. RMN: \ce{OCH2}
4.13 (c), \ce{CH2CO} 2.28 (t), \ce{CH2} central 1.66 (sextuplete),
\ce{CH3} del éster 1.25 (t), \ce{CH3} de la cadena 0.95 (t),
\omterm{def:b1:structure-spectroscopy:integration}{integraciones} 2, 2, 2, 3, 3.
\textbf{26.} \textbf{\qty{116}{g/mol}}: el butanoato de etilo,
\ce{C6H12O2}.
\end{solution}
